% ===== PRÉAMBULE CANONIQUE — à coller VERBATIM en tête de CHAQUE .tex =====
\documentclass[11pt,a4paper]{article}
\usepackage[utf8]{inputenc}\usepackage[T1]{fontenc}\usepackage{lmodern}
\usepackage[french]{babel}
\usepackage{amsmath,amssymb,mathtools}\usepackage{icomma}
\usepackage{siunitx}\sisetup{locale=FR}  % \qty{}{}, \si{}, \ang{}
\usepackage{xcolor}\usepackage{colortbl}
\usepackage{tikz}\usepackage{pgfplots}\pgfplotsset{compat=1.18}
\usepackage[siunitx,europeanresistors]{circuitikz} % circuits (élec.)
\usepackage[most]{tcolorbox}\usepackage{enumitem}\usepackage{array,booktabs}
\usepackage[a4paper,margin=2.2cm]{geometry}
\usetikzlibrary{arrows.meta,calc,angles,quotes,positioning,patterns,%
  backgrounds,shapes.geometric,decorations.pathmorphing,decorations.markings,intersections}
\definecolor{primary}{RGB}{222,49,99}\definecolor{primarydark}{RGB}{150,22,55}
\definecolor{defcol}{RGB}{0,114,178}\definecolor{methcol}{RGB}{52,92,148}
\definecolor{excol}{RGB}{0,135,90}\definecolor{attncol}{RGB}{200,40,55}
\tcbset{boxbase/.style={enhanced,breakable,boxrule=0.9pt,arc=2.5pt,
  left=3mm,right=3mm,top=2.3mm,bottom=2.3mm,fonttitle=\bfseries,coltitle=white}}
% --- boîtes TUTORIEL (noms EXACTS, ne pas renommer) ---
\newtcolorbox{cle}[1][]{boxbase,colback=defcol!6,colframe=defcol,
  title={\textsc{À retenir}\ifx\relax#1\relax\relax\else\textmd{ : \itshape #1}\fi}}
\newtcolorbox{methode}[1][]{boxbase,colback=methcol!7,colframe=methcol,sharp corners,
  title={\textsc{Méthode}\ifx\relax#1\relax\relax\else\textmd{ --- \itshape #1}\fi}}
\newtcolorbox{exemplebox}[1][]{boxbase,colback=excol!5,colframe=excol,
  title={\textsc{Exemple}\ifx\relax#1\relax\relax\else\textmd{ : \itshape #1}\fi}}
\newtcolorbox{piege}[1][]{boxbase,colback=attncol!6,colframe=attncol,
  title={\textsc{Piège}\ifx\relax#1\relax\relax\else\textmd{ : \itshape #1}\fi}}
\newtcolorbox{remarque}[1][]{boxbase,colback=black!4,colframe=black!45,
  title={\textsc{Remarque}\ifx\relax#1\relax\relax\else\textmd{ : \itshape #1}\fi}}
% --- boîtes EXERCICE / CORRIGÉ (noms EXACTS, auto-comptées) ---
\newtcolorbox[auto counter]{exo}[1][]{boxbase,colback=primary!4,colframe=primary,
  title={\textsc{Exercice}~\thetcbcounter\ifx\relax#1\relax\relax\else\textmd{ --- \itshape #1}\fi}}
\newtcolorbox[auto counter]{corrige}[1][]{boxbase,colback=excol!4,colframe=excol,
  title={\textsc{Corrigé}~\thetcbcounter\ifx\relax#1\relax\relax\else\textmd{ --- \itshape #1}\fi}}
\newcommand{\vv}[1]{\vec{#1}}\newcommand{\ds}{\displaystyle}
\newcommand{\R}{\mathbb{R}}\newcommand{\N}{\mathbb{N}}\newcommand{\Z}{\mathbb{Z}}
% =========================================================================
\begin{document}
\usetikzlibrary{babel}
\begin{corrige}[énergie de Joule dégagée]
$t=\qty{5}{\minute}=300\,\si{\second}$. Loi de Joule :
\[ W_{\text{th}}=R\,I^{2}\,t=20\times3^{2}\times300=20\times9\times300=\qty{54000}{\joule}=\qty{54}{\kilo\joule}. \]
\end{corrige}

\begin{corrige}[plaque signalétique d'un appareil]
\begin{enumerate}[label=\alph*)]
  \item $I=\dfrac{P}{U}=\dfrac{2000}{230}\approx\qty{8.7}{\ampere}$.
  \item $R=\dfrac{U^{2}}{P}=\dfrac{230^{2}}{2000}=\dfrac{52900}{2000}\approx\qty{26.5}{\ohm}$.
\end{enumerate}
\end{corrige}

\begin{corrige}[trois écritures, un seul résultat]
\begin{enumerate}[label=\alph*)]
  \item $I=\dfrac{U}{R}=\dfrac{24}{12}=\qty{2}{\ampere}$.
  \item $P=U\,I=24\times2=\qty{48}{\watt}$ ; \; $P=R\,I^{2}=12\times2^{2}=\qty{48}{\watt}$ ; \;
        $P=\dfrac{U^{2}}{R}=\dfrac{576}{12}=\qty{48}{\watt}$. Les trois concordent \checkmark.
\end{enumerate}
\end{corrige}

\begin{corrige}[énergie et facture]
\begin{enumerate}[label=\alph*)]
  \item $W=P\,t=2\times3=\qty{6}{\kilo\watt\hour}$ par jour.
  \item Coût $=6\times0{,}15=\num{0.90}$ euro par jour.
\end{enumerate}
\end{corrige}

\begin{corrige}[à même tension, qui chauffe le plus ?]
\begin{enumerate}[label=\alph*)]
  \item $P_1=\dfrac{U^{2}}{R_1}=\dfrac{400}{10}=\qty{40}{\watt}$ ; \;
        $P_2=\dfrac{U^{2}}{R_2}=\dfrac{400}{40}=\qty{10}{\watt}$.
  \item À \emph{même tension}, $P=U^2/R$ : c'est la \textbf{plus petite résistance} qui dissipe le plus
        ($R_1$ chauffe $4$ fois plus que $R_2$). Elle laisse passer plus de courant.
\end{enumerate}
\end{corrige}

\end{document}
